% 04e-lp.tex — v3 §3.6 (first part): Lemma 3.12 and Theorem 3.13 (local PNT) with full proof.
\subsection{A local prime number theorem}\label{sec:noise-lp}
% v3: §3.6 (Lemma 3.12, Theorem 3.13)

\begin{lemma}[threshold mass]\label{lem:threshold-mass}
% v3: Lemma 3.12
For squarefree $n$ put $\tau(n):=\max_{p\mid n}t(p)\prod_{q\mid n,\,t(q)\ge t(p)}q$, so that $\chi^t_y(n)=1$ if and only if
$\tau(n)<y$. Assume $|N-T|\le D:=B+\frac12H(1+\kappa_+/\log x)+4$ at the relevant heights (Corollary~\ref{cor:lag}). For $x\ge X_0^3$ and
% Logged correction 38 (CORRECTIONS-vs-v3.md): x >= X0^3 (v3: x >= X0), so that every p in case (ii) exceeds X0.
% Logged correction 36 (CORRECTIONS-vs-v3.md): +4 (lower lag constant -4); D <= H still, as B << H.
$\delta\in(0,\frac14]$,
\[
  \sum_{n:\ \tau(n)\in[x(1-\delta),x)}\frac1n\le55\,\delta\log x\,\big(\log H(x)+\log(1/\delta)+2\big)
  +14\big(K_{\mathrm{loc}}(\sqrt x)+2\big)\log x\,\sqrt{\frac{H(x)}{\delta x}} .
\]
It is used only with $\delta\ge6H(x)x^{-1/3}$, and only in Proposition~\ref{prop:g1}(B), where $x\ge X_{\mathrm{sw}}$ with
$\log X_{\mathrm{sw}}=X_0^{1/100}>3\Lambda$.
\end{lemma}

\begin{proof}
Let $p$ be the prime attaining $\tau(n)=t(p)\,m$, with $m:=\prod_{t(q)\ge t(p)}q$, and write $n=k\cdot m$, where all
primes of $k$ are $t$-below $p$. The map $n\mapsto(k,m)$ is injective. A prime $q$ with $t(q)<t(p)$ has $q<P:=p+3H(p)$,
and $q<p$ if $p<\Xm$ (then $q$ and $p$ are rational primes, with $t(q)=q$). Hence
$\sum_k1/k\le\prod_{q<P}(1-1/q)^{-1}$, which we bound as follows.
\begin{itemize}
\item If $p<\Xm$: for $p<400$ a direct computation gives $\prod_{q<p}(1-1/q)^{-1}\le1.8\log p+0.25$ (Appendix~\ref{app:cert}, \texttt{finite-n}); for $p\ge400$,
  Rosser and Schoenfeld \cite[Thm.~8, (3.29)]{RosserSchoenfeld1962} give
  $\prod_{q\le p}(1-1/q)^{-1}\le e^\gamma\log p\,(1+1/(2\log^2p))\le1.7811\log p+0.16$.
\item If $p\ge\Xm$: split at $y_0:=\Xm-1$, below which all primes are rational. The rational part is at most
  $e^\gamma\log y_0\,(1+1/(2\log^2y_0))$ by the same bound. For the rest, $-\log(1-1/q)\le1/q+1/q^2$, and
  Lemma~\ref{lem:cm}(iv) with $y=y_0$, $x=P$ gives $\sum_{y_0<q<P}1/q\le\log(\log P/\log y_0)+2.2/\log y_0$, while
  $\sum_{q>y_0}1/q^2\le3/y_0$. So the product is at most $1.7811\log P\,(1+3\cdot10^{-8})\le1.782\log p$.
\end{itemize}
In all cases $\sum_k1/k\le1.8\log p+1\le2\log x$, as $p\le x$ and $\log x\ge10^8$. Also $1/m\le2p/x$.

Write $m=p\,m'$, so that $p^2m'\in I:=[x(1-\delta),x]\,(1+O(H(p)/p))$. In case (ii) below, $H(p)/p\le\delta/10$, since
$\delta\ge6H(x)x^{-1/3}$.
\begin{enumerate}[label=(\roman*)]
\item $p\le\sqrt{\delta x/(2H(x))}$. The interval $I/p^2$ has length $\ge2D$, so it holds at most $4\delta x/p^2$
  elements $m'$ (by $|N-T|\le D$). These pairs contribute at most $4\delta x\sum_{p\le\sqrt x}1/p$ to $\sum p$.
\item $p$ larger. Then $m'\le2H(x)/\delta$. For each $m'$, the $p$-interval $\sqrt{I/m'}$ has length
  $\ell_{m'}\le0.6\,\delta\sqrt{x/m'}$. Since $p^2m'=pm=\tau(n)\,p/t(p)$ and $|t(p)-p|\le H(p)\le\delta p/10$, the factor
  $1+O(H(p)/p)$ in $I$ lies within $1\pm0.11\delta$. So, with $m'\le2H(x)/\delta$, $\delta\le\frac14$, $\delta\ge6H(x)x^{-1/3}$ and
  $x\ge X_0^3$, the interval has minimum $\ge\sqrt{(1-\delta)(1-0.11\delta)}\,\sqrt{\delta x/(2H(x))}\ge0.85\sqrt3\,x^{1/3}\ge1.47X_0$
  and endpoint ratio $\le\big((1+0.11\delta)/((1-\delta)(1-0.11\delta))\big)^{1/2}\le1.19$. Let $P_1\le\sqrt x$ be the dyadic scale
  with $P_1\le\min\sqrt{I/m'}<2P_1$; then $P_1\ge X_0$ and the interval lies in $[P_1,2.38P_1)\subset[P_1,4P_1)$. Cut it into
  $\lceil\ell_{m'}/(6H(P_1))\rceil$ windows of length $\le6H(P_1)$ and apply the local count at the dyadic scale $P_1$
  (Corollary~\ref{cor:localcount}(i)) to each. This gives at most
  $13.4\ell_{m'}+\lceil\ell_{m'}/(6H(P_1))\rceil(K_{\mathrm{loc}}(P_1)+2)\le13.5\ell_{m'}+K_{\mathrm{loc}}(\sqrt x)+2$ primes, since
  $K_{\mathrm{loc}}\le H/100$ and $K_{\mathrm{loc}}$ is non-decreasing.
% Logged correction 37 (CORRECTIONS-vs-v3.md): v3 applies the local count at the non-dyadic height sqrt(x). Each p-interval
% lies in one block [P_1,4P_1) with P_1 dyadic, so part (i) applies there and the constant K_loc(sqrt x) is unchanged.
% (Part (ii), for real P, would give K_loc(P)+K_loc(2P) and turn the constant 14 of the statement into 23.)
\end{enumerate}
Each pair $(p,m')$ contributes at most $2\log x\cdot2p/x$. Collecting, with $p\le\sqrt{x/m'}$, the sum is at most
\[
  \frac{4\log x}x\Big[4\delta x(\log\log x+1)+8.1\,\delta x\big(\log(2H(x)/\delta)+1\big)+2(K_{\mathrm{loc}}+2)\sqrt x\,\sqrt{2H(x)/\delta}\Big].
\]
Since $\log H\ge8(\log\log x)^2\ge16\log\log x$, this is at most the stated bound.
\end{proof}

\paragraph{Notation for the local prime number theorem.}
\begin{itemize}
\item $\Pi$ is the realised prime measure (prime powers weighted $1/k$), so that $N=\exp^*(\Pi)$; and
  $\Pi_0:=\log^*(\delta_1+du)$, whose density is $(1-1/u)/\log u$.
\item $T_0:=T-(\#\Abs-\int\bar\rho_{\Abs})$ is the $\Abs$-free target. Its Mellin transform is
  $\hat T_0(s)=s/(s-1)+\hat R(s)+\hat E(s)$ (Lemma~\ref{lem:ZBlow}), where $\hat R=\sum_k\hat R_k$ is the transform of the
  designed oscillation \eqref{eq:Rhat}.
\item $\theta':=\max(\theta+0.01,\ 8(\log\log X_0)^2/\log X_0+0.005)\le0.51$ is a fixed exponent (only $\theta'<1$ is
  needed), and $\Lambda=\log X_0$.
\end{itemize}

\begin{theorem}[local prime number theorem]\label{thm:lp}
% v3: Theorem 3.13
Let $y\ge X_0$, and assume:
\begin{enumerate}[label=\textup{(LP-\alph*)}]
\item $N=T_0$ below $\Xm$, and $|N(u)-T_0(u)|\le D_0(u):=B(u)+3H(u)+4+3.1u^{\theta/2}(\log u)^{-4}\le u^{\theta'}$ for
  $\Xm\le u\le y$. (On $\Good(<y)$ this holds by Theorem~\ref{thm:queue}, Corollary~\ref{cor:lag} and \eqref{eq:T1}.)
% Logged correction 34 (CORRECTIONS-vs-v3.md): 3.1 (v3: 2) in D_0, for the accumulated mass deficit of the smoothed
% absorbed elements. D_0 <= u^{theta'} keeps huge room: theta' >= 0.26, and 3H(u) <= 3exp(8(loglog u)^2 + 1e-6 log u).
\item $|\hat E(s)|\le C_8(1+\kappa|s|)e^\gamma\Lambda$ for $\sigma\ge1-8/\Lambda$, with $\log C_8<440$
  (Lemma~\ref{lem:ZBlow} with $a=8$).
\item $|\hat R(s)|\le|s|$ for $\sigma\ge1-8/\Lambda$. (By \eqref{eq:Rhat}: there $\mathrm{Re}\,w=\sigma-\theta\ge0.4$, so
  $|\hat R_k(s)|\le C\alpha_k|s|x_k^{\theta-\sigma}\le2C|s|x_k^{-0.4}$, which is summable over $k$ since $x_1\ge\exp(X_0^2)$;
  see also \eqref{eq:Rhatbound}.)
\end{enumerate}
Then there are $c_{\mathrm{zf}}\ge0.088/\Lambda$, $c_{\mathrm{LP}}:=c_{\mathrm{zf}}/5\ge0.0176/\Lambda$, an effective $K$ with
$\log K=O(\Lambda)$, and a threshold $X_{\mathrm{LP}}$ with $\log X_{\mathrm{LP}}=O(\Lambda^2\log^2\Lambda)$, such that for
$X_{\mathrm{LP}}\le v\le y$:
\begin{enumerate}[label=\textup{(\roman*)}]
\item $\displaystyle\Big|\int_{[1,v]}\frac{d(\Pi-\Pi_0)(u)}u\Big|\le K\,e^{-c_{\mathrm{LP}}\sqrt{\log v}}$;
\item $|\Pi(v)-\Pi_0(v)|\le K\,v\,e^{-c_{\mathrm{LP}}\sqrt{\log v}}$;
\item \emph{(short windows)} let $\Lambda_p:=\log v\ge76\Lambda^3$, $\eta\in[e^{-2\Lambda},\frac12]$, and let $\varphi\ge\mathbf 1_{[v,v(1+\eta)]}$ be the
  smooth window of relative width $\eta(1+2\Lambda_p^{-2})$ with transition layers $\eta\Lambda_p^{-2}$. Put
  $L_1:=2\log\Lambda_p+\log(1/\eta)+5$. Then $\int\varphi\,d\Pi\le(1+\mathrm{Rel})\int\varphi\,d\Pi_0$, with
  \begin{equation}\label{eq:lpplus}
    \mathrm{Rel}(\Lambda_p,\eta)\le18\,G(L_1)\,\Lambda_p^3\,\eta^{-1}\exp\Big(-\frac{0.9\,c_{\mathrm{zf}}\Lambda_p}{L_1+442+2\log\Lambda+\log L_1}\Big),
  \end{equation}
  where $G(L):=\sup|\log(\zeta^\sharp(s)(s-1)/s)|$ on the shifted contour up to height $e^L$. Moreover
  $G(L)\le400(L+450+2\log\Lambda)$ for $5\le L\le\log T_y$, with $\log T_y:=\log y/(51\Lambda\log\log y)-500$; so
  $\log G(L)\le\Lambda$ for such $L$ with $L\le e^{\Lambda/2}$.
% Logged correction 24 (CORRECTIONS-vs-v3.md): (iii) restricted to Lambda_p >= 76 Lambda^3; G defined via the integrand
% log(zeta#(s)(s-1)/s) with the explicit bound above (v3: L_1 <= 3000 Lambda and log G <= Lambda unconditionally).
\end{enumerate}
\end{theorem}

\begin{proof}[Proof of Theorem~\ref{thm:lp}]
\begingroup\emergencystretch=3em
This is Landau's argument \cite[Lemmas~6.2--6.3 and the proof of Thm.~6.6]{MontgomeryVaughan2007}, with four changes. It is
run on an extension of $N$ beyond $y$, which is not a Beurling system. Positivity is therefore imported from a finite
Euler product. Landau's discs have radius $\asymp1/\Lambda$. And the region near the pole is handled with the fine form of
$\hat E$. Throughout, $s=\sigma+it$.\par\endgroup

\emph{Step 1: the extension.} Put $D:=N-T_0$, so $D=0$ below $\Xm$ and $|D(u)|\le u^{\theta'}$ on $[\Xm,y]$ by (LP-a). Define
$N^\sharp:=N$ on $[1,y]$ and $N^\sharp(u):=N(y)+T_0(u)-T_0(y)$ for $u>y$. Then $D^\sharp:=N^\sharp-T_0$ equals
$D(\min(u,y))$, and
\[
  \zeta^\sharp(s):=\int_{[1,\infty)}u^{-s}\,dN^\sharp=\hat T_0(s)+\hat D^\sharp(s),\qquad
  \hat D^\sharp(s)=s\int_{\Xm}^yD(u)u^{-s-1}du+D(y)y^{-s}.
\]
So $\hat D^\sharp$ is entire, with $|\hat D^\sharp(s)|\le|s|(\Xm)^{\theta'-\sigma}/(\sigma-\theta')+y^{\theta'-\sigma}$ for
$\sigma>\theta'$. Let $\Pi^\sharp:=\log^*(dN^\sharp)$, a signed measure (the series converges in the weighted norm
$\int u^{-\sigma}|\cdot|$ for large $\sigma$). Since $\log^*$ is causal, $\Pi^\sharp=\Pi$ on $[1,y]$.

\emph{Step 2: upper bounds on $[1,y]$.} Lemma~\ref{lem:cm} holds on $[1,y]$: (LP-a) and \eqref{eq:T1} give \textup{(H$_{\mathrm{BS}}$)} there. Write $d\psi=\log u\,d\Pi$ and
$S(z)=\int_{[1,z]}d\psi/u\le\log z+0.43$. By Lemma~\ref{lem:cm}(iv) and the values for $\N$,
$\int_{[1,y]}d\Pi/u\le\log\log y+4.5$. Let $N_y:=\exp^*(\Pi|_{[1,y]})$ count the $y$-smooth elements, so $N_y=N$ on
$[1,y]$, and put $P_y(s):=\int u^{-s}dN_y=\exp(\int_{[1,y]}u^{-s}d\Pi)$. Then $P_y(1)\le e^{4.5}\log y$.

\emph{Step 3: positivity transfer.} Put $\Delta:=\zeta^\sharp-P_y=\int_{(y,\infty)}u^{-s}(dT_0-dN_y)$. For $\sigma=1+\eta$
with $0<\eta\le1/8$, Rankin's trick ($u^{-\eta}\le y^{-\eta/2}u^{-\eta/2}$ for $u>y$) and (LP-a) give
\begin{gather*}
  P_y(\sigma)\le1/\eta+6+y^{-\eta/2}P_y(1),\qquad|\Delta(s)|\le2y^{-\eta}/\eta+y^{-\eta/2}P_y(1),\\
  |\Delta'(s)|\le2y^{-\eta}(\log y/\eta+1/\eta^2)+y^{-\eta/2}P_y(1)\,S(y),\qquad|P_y(s)|\ge1/P_y(\sigma),\qquad
  |P_y'/P_y(s)|\le S(y).
\end{gather*}
The last two hold because $\Pi\ge0$. Set $\eta_y:=17\log\log y/\log y$, so that $y^{-\eta/2}\le(\log y)^{-8.5}$ for
$\eta\ge\eta_y$. Then, for $1+\eta_y\le\sigma\le1+1/8$,
\begin{equation}\label{eq:lp-transfer}
  P_y(\sigma)\le2/\eta,\qquad|\Delta/P_y|\le10^{-2},\qquad\Big|\frac{d}{ds}\log(1+\Delta/P_y)\Big|\le(\eta\log y)^{-1}.
\end{equation}
Let $F:=(dT_0-dN_y)|_{(y,\infty)}*(dN_y)^{*-1}$. Since $|(dN_y)^{*-1}|\le dN_y$, the same bounds give
$\int u^{-\sigma}|dF|\le10^{-2}$. Hence $\Pi^\sharp=\Pi|_{[1,y]}+\log^*(\delta_1+F)$ satisfies $\int u^{-\sigma}|d\Pi^\sharp|<\infty$,
$\zeta^\sharp$ has no zeros in $\sigma\ge1+\eta_y$, and $\log\zeta^\sharp=\log P_y+\log(1+\Delta/P_y)$ is the Mellin transform of
$\Pi^\sharp$ there. The coefficients of $-P_y'/P_y$ form the positive measure $d\psi|_{[1,y]}$. So for the nonnegative
polynomial
\[
  (1+\cos\phi)^4=\tfrac18\big(35+56\cos\phi+28\cos2\phi+8\cos3\phi+\cos4\phi\big),
\]
with $(a_0,\dots,a_4):=(35,56,28,8,1)$, we get
\begin{equation}\label{eq:lp-341}
  \sum_{k=0}^4a_k\,\mathrm{Re}\Big(-\frac{\zeta^{\sharp\prime}}{\zeta^\sharp}\Big)(\sigma+ikt)\ge-\frac{128\,\eps_f}{\eta},
  \qquad\eps_f:=5\cdot10^{-4}\ \ (\ge1/\log y),\qquad\sigma=1+\eta\ge1+\eta_y .
\end{equation}

\emph{Step 4: the pole term.} By Step~3 and Lemma~\ref{lem:cm}(ii), with partial summation,
$-P_y'/P_y(1+\eta)=\int_{[1,y]}u^{-\eta}dS(u)\le\eta\int_1^\infty(\log u+0.43)u^{-1-\eta}du=1/\eta+0.43$. Hence, by
\eqref{eq:lp-transfer},
\begin{equation}\label{eq:lp-pole}
  -\frac{\zeta^{\sharp\prime}}{\zeta^\sharp}(1+\eta)\le\frac{1+10^{-8}}\eta+0.43\qquad(\eta\ge\eta_y).
\end{equation}
This replaces the usual use of the pole of $\zeta^\sharp$; the crude bound $|\hat E|\le C_8e^\gamma\Lambda$ is useless near
$s=1$.

\emph{Step 5: Landau's lemma.} Let $f$ be analytic on $|z-z_0|\le r$, with $|f|\le M|f(z_0)|$ there, and suppose every zero
of $f$ in the disc has real part $\le\mathrm{Re}\,z_0$. Then for any zero $\rho_1$ in the disc,
\[
  -\mathrm{Re}\frac{f'}f(z_0)\le\frac2r\log M-\mathrm{Re}(z_0-\rho_1)\Big(\frac1{|z_0-\rho_1|^2}-\frac1{r^2}\Big).
\]
\emph{Proof.} Divide $f$ by the Blaschke factors of its zeros, normalised on the circle of radius $r$; the quotient $g$ has
no zeros and $|g/g(z_0)|\le M$ on the circle. Borel--Carath\'eodory for $h=\log(g/g(z_0))$ on radii $R\uparrow r$ gives
$|h'(z_0)|\le2\log M/r$. Each zero contributes $\mathrm{Re}(z_0-\rho)(1/|z_0-\rho|^2-1/r^2)\ge0$ to $-\mathrm{Re}\,f'/f(z_0)$;
drop all but $\rho_1$.

\emph{Step 6: the region near the pole.} For $|t|\le X_0$ and $\sigma\ge1-8/\Lambda$, the proof of
Lemma~\ref{lem:ZBlow} gives the fine form
$\hat E(s)=\zeta(s)-1/(s-1)-1+O(X_0^{-\sigma})+O(\kappa(1+|s|)e^{\mathrm{Ein}(8)}e^\gamma\Lambda)$. Since
$s/(s-1)=1+1/(s-1)$, \eqref{eq:Rhatbound} (that is, $|\hat R(s)|\le|s|e^{-0.39X_0^2}$) and Step~1 give
\[
  \zeta^\sharp(s)=\zeta(s)+E(s),\qquad|E(s)|\le3X_0^{-0.3}(1+|t|)\qquad(|t|\le X_0,\ \sigma\ge1-8/\Lambda),
\]
using $\kappa\le X_0^{-1/3}$ and $e^{\mathrm{Ein}(8)}\le e^{438}$. On $|s-1|\le1/8$ the Laurent series gives
$|\zeta(s)-1/(s-1)|\le\gamma+0.02\le0.6$ (as $|\gamma_1|\le0.073$), so $|(s-1)\zeta^\sharp(s)-1|\le0.6/8+10^{-9}<1$, and
$\zeta^\sharp$ has no zero there. This covers every point with $|t|<1/16$ and $1-8/\Lambda\le\sigma\le1.1$; for
$\sigma>1.1$, $\zeta^\sharp\ne0$ by Step~3.

\emph{Step 7: the zero-free region for $|t|\ge1/16$.} Take Landau discs of radius $r:=a/\Lambda$ with $a:=6.8\le8$, so that
they stay in $\sigma\ge1-8/\Lambda$, where (LP-b) and (LP-c) apply. For $|t|\ge1/16$ they do not meet the pole. Put
\[
  L_t:=\log(|t|+3),\qquad X(t):=L_t+442+2\log\Lambda,\qquad Y(t):=X(t)+\log X(t)+0.7 .
\]
\begin{itemize}
\item \emph{Growth.} Let $|t|\ge1/16$, $\eta_y\le\eta\le10^{-3}/\Lambda$, $k\in\{1,2,3,4\}$ and $z_0=1+\eta+ikt$. Step~3 gives
  $|\zeta^\sharp(z_0)|\ge(1-10^{-2})/P_y(1+\eta)\ge0.495\,\eta$. On the disc $|z-z_0|\le r$:
  \begin{itemize}
  \item $|z/(z-1)|\le1+1/(k|t|-r)\le17.1$;
  \item $|\hat R(z)|\le|z|\le k|t|+1.01$, by (LP-c);
  \item $|\hat E(z)|\le e^\gamma C_8\Lambda+10^{-10^7}|z|$, by (LP-b), since $\kappa\le X_0^{-1/3}$ and $\log C_8<440$;
  \item $|\hat D^\sharp(z)|\le2.1|z|X_0^{-0.48}+1$, by Step~1, since $\sigma-\theta'\ge0.48$.
  \end{itemize}
  So $|\zeta^\sharp|\le2(1+3X_0^{-0.48})\max(k|t|+19.2,\ e^\gamma C_8\Lambda)$ there, for every $|t|\ge1/16$. Now
  $\log(4|t|+19.2)\le L_t+1.857$ and $\log(e^\gamma C_8\Lambda)<440.578+\log\Lambda$; since $\max(a,b)\le a+b-2.975$ when
  $a,b\ge2.975$, the ratio $M_k:=\max_{|z-z_0|\le r}|\zeta^\sharp(z)|/|\zeta^\sharp(z_0)|$ satisfies (the constant is
  $440.855+\log(1+3X_0^{-0.48})\le440.86$)
  \begin{equation}\label{eq:lp-growth}
    \log M_k\le L_t+440.86+\log\Lambda+\log(1/\eta).
  \end{equation}
\item \emph{The inequality.} Suppose $\zeta^\sharp(\beta+it)=0$ with $1-\beta\le\eta/\kappa_L$ for a parameter $\kappa_L>0$, and
  $\eta\ge\eta_y$. Apply Step~5 at $z_0=1+\eta+it$ with this zero, and at $z_0=1+\eta+ikt$ ($k=2,3,4$) with no zero term.
  This is legitimate because all zeros have real part $<1+\eta_y\le1+\eta$ (Step~3), so the dropped terms are $\ge0$. Insert
  these and \eqref{eq:lp-pole} into \eqref{eq:lp-341}. Using $1/(\eta+1-\beta)\ge\kappa_L/((\kappa_L+1)\eta)$, we get
  \[
    \frac{N_0(\kappa_L)}{\eta}\le15.05+\frac2r\sum_{k=1}^4a_k\log M_k+\frac{56\,\eta(1+1/\kappa_L)}{r^2},\qquad
    N_0(\kappa_L):=\frac{56\kappa_L}{\kappa_L+1}-35(1+10^{-8})-128\eps_f .
  \]
\item \emph{The constant.} Since $\sum_{k=1}^4a_k=93$, the inequality gives, as $|t|\to\infty$, the zero-free width
  $\mathrm{merit}\cdot r/(2\log|t|)$, with
  \[
    \mathrm{merit}:=\max_{\kappa_L}\frac{N_0(\kappa_L)}{93\,\kappa_L}=0.0262291\qquad(\kappa_L\approx3.79).
  \]
  So $\mathrm{merit}\cdot r/2=0.0891790/\Lambda$, and $\mathrm{merit}\cdot a/2\ge0.088$ needs $a\ge6.71$. We fix $\kappa_L:=3.79$,
  for which $N_0\ge9.2449$, and put
  \[
    c_{\mathrm{zf}}:=0.088/\Lambda .
  \]
  The slack between $0.08917$ and $0.088$ absorbs the lower-order terms of \eqref{eq:lp-growth} at finite height, as follows.
  The constant $\mathrm{merit}$ does not depend on $\Lambda$.
\item \emph{Explicit width.} Let $|t|\ge1/16$, and suppose that $\eta:=0.33769/(\Lambda Y(t))$ satisfies $\eta\ge\eta_y$. Then
  $\zeta^\sharp(\sigma+it)\ne0$ for
  \begin{equation}\label{eq:lp-width}
    \sigma\ge1-\delta_t,\qquad\delta_t:=\eta/\kappa_L=\frac{0.0891}{\Lambda\,Y(t)} .
  \end{equation}
  Indeed, otherwise the inequality holds with this $\eta$. Here $\log(1/\eta)=\log\Lambda+\log Y+1.0857$, and
  $\log Y-\log X\le(\log X+0.7)/X\le0.0144$ (as $X\ge479.96$), so the right side of \eqref{eq:lp-growth} is at most
  $X-0.054+\log Y\le X+\log X-0.039\le Y$.
  Multiply the inequality by $\eta$, and use $\eta\Lambda=0.33769/Y\le7.1\cdot10^{-4}$: the right side becomes at most
  $(2\cdot93/6.8)\cdot0.33769+10^{-6}\le9.2369<9.2449\le N_0$, a contradiction.
\item \emph{Two consequences.} Write $\delta_{\mathrm{zf}}(t):=\delta_t$.
  \begin{itemize}
  \item Since $0.088(X+\log X+0.7)\le0.0891(X+18)$ for $X\ge479$, we have $\delta_{\mathrm{zf}}(t)\ge c_{\mathrm{zf}}/(L_t+460+2\log\Lambda)$.
  \item \emph{Finite height.} Let $L\ge5$ and $|t|\le e^L$. Then $L_t\le L+3e^{-L}$, and
    $0.088(X+\log X+0.7)\le0.0891(L+442+2\log\Lambda+\log L)$: the difference of the two sides increases with $L$ and with
    $\Lambda$, and is positive at $L=5$, $\Lambda=10^8$. So at height $e^L$ the zero-free width is
    $\delta_{\mathrm{zf}}\ge c_{\mathrm{zf}}/(L+442+2\log\Lambda+\log L)$. Every use below has $L\ge40$.
  \end{itemize}
\item \emph{The height constraint.} Step~3 requires $\eta=\kappa_L\delta_{\mathrm{zf}}(t)\ge\eta_y=17\log\log y/\log y$, and under the
  present scaling this binds. Assume
  \begin{equation}\label{eq:lp-yrange}
    \log y\ge10^4\Lambda^2\log\log y,\qquad\text{so that}\quad Q_y:=\frac{\log y}{\Lambda\log\log y}\ge10^4\Lambda ,
  \end{equation}
  and put $\log T_y:=Q_y/51-500$. For $|t|\le T_y$ we have $X(t)\le Q_y$ and
  $Y(t)\le Q_y/51-57.3+2\log\Lambda+\log Q_y\le Q_y/50.35$, because $(1/50.35-1/51)Q_y\ge2.53\cdot10^{-4}Q_y\ge2\log\Lambda+\log Q_y$.
  Since $0.33769\cdot50.35\ge17.002$, this gives $\eta=0.33769/(\Lambda Y)\ge17/(\Lambda Q_y)=\eta_y$. So the zero-free region is
  \begin{equation}\label{eq:lp-zf}
    \sigma\ge1-\frac{c_{\mathrm{zf}}}{\log(|t|+3)+460+2\log\Lambda}\qquad\text{for }|t|\le T_y,\qquad
    \log T_y:=\frac{\log y}{51\Lambda\log\log y}-500 .
  \end{equation}
  Together with Step~6 this covers all $|t|\le T_y$. Every application below satisfies \eqref{eq:lp-yrange}.
\item By (LP-c), $\hat R$ enters only through the growth bound, like $\hat D^\sharp$. The region \eqref{eq:lp-zf} is
  consistent with the designed zeros, which have $1-\beta_k\asymp\log\log\gamma_k/\log\gamma_k$ (\S\ref{sec:zeros}).
\end{itemize}

\emph{Step 8: $\log\zeta^\sharp$ in the region.} Assume \eqref{eq:lp-yrange}, and let $\mathfrak G_1(s):=\log(\zeta^\sharp(s)(s-1)/s)$,
the branch that is real on $(1,\infty)$; it is analytic in the zero-free region, including near $s=1$. Fix $f:=0.9$, and
for $1/16\le|t|\le T_y$ let $s_1:=1+\eta_t+it$ with $\eta_t:=\kappa_L\delta_t$. Consider the discs centred at $s_1$ of
radii $R_1=\eta_t+\frac{1+f}2\delta_t$ and $r_1=\eta_t+f\delta_t$. They avoid the pole, and they lie in the zero-free
region, since $\delta_{\mathrm{zf}}$ changes by a relative $10^{-8}$ over them. The function
$h:=\log\zeta^\sharp-\log\zeta^\sharp(s_1)$ vanishes at $s_1$, and $\mathrm{Re}\,h\le\log M_1\le Y(t)$ on the larger disc, by
\eqref{eq:lp-growth} with $\eta=\eta_t$ (as in Step~7). Borel--Carath\'eodory gives
$|h|\le\frac{2r_1}{R_1-r_1}Y(t)=\frac{4(\kappa_L+f)}{1-f}Y(t)\le187.6\,Y(t)$ on the smaller disc. At the anchor, Step~3 gives
$|\log\zeta^\sharp(s_1)|\le\log P_y(1+\eta_t)+0.02\le\log(2/\eta_t)+0.02\le0.5\,Y(t)$. Finally
$|\log((s-1)/s)|\le\log17+\pi\le6$ for $|t|\ge1/16$ and $\sigma\le1.1$. Hence
\[
  |\mathfrak G_1(s)|\le190\,Y(t)\qquad\text{for }|s-s_1|\le r_1 .
\]
A point $s=\sigma+it$ of the contour $\sigma=1-f\delta_{\mathrm{zf}}(T_1)$, with $1/16\le|t|\le T_1\le T_y$, satisfies
$|s-s_1|=\eta_t+f\delta_{\mathrm{zf}}(T_1)\le r_1$, because $\delta_{\mathrm{zf}}$ decreases in $|t|$. For $|t|<1/16$ the contour
lies in $|s-1|\le1/16+10^{-8}$. There Step~6 gives $|\log((s-1)\zeta^\sharp(s))|\le-\log(1-0.075)\le0.078$, and
$|\log s|\le-\log(1-|s-1|)\le0.0646$, so $|\mathfrak G_1|\le0.1426$. Since
$Y\le2X$, it follows that for $5\le L\le\log T_y$ the quantity $G(L)$ of (iii) satisfies
\[
  G(L)\le190\,Y(e^L)\le400\,(L+450+2\log\Lambda);
\]
in particular $\log G(L)\le\Lambda$ for $L\le e^{\Lambda/2}$.

\emph{Step 9: Perron's formula; proof of (i) and (ii).} Let
$\mathfrak G(s):=\log\big(\zeta^\sharp(s+1)\,s/(s+1)\big)$. It is analytic in the shifted zero-free region, $\mathfrak G(0)=0$
(the residue of $\zeta^\sharp$ at $1$ is $1$, and $\zeta^\sharp>0$ on $(1,\infty)$), and $\mathfrak G$ is the Mellin transform of
$(\Pi^\sharp-\Pi_0)/u$ on $\mathrm{Re}\,s\ge\eta_y$. Let $\Phi^\sharp(w):=\int_{[1,w]}d(\Pi^\sharp-\Pi_0)/u$; it equals the left
side of (i) for $w\le y$.
\begin{itemize}
\item Take $h_0\in C_c^\infty(-2,-1)$ with $h_0\ge0$, $\int h_0=1$, and put $h:=Th_0(T\cdot)$. Smoothing to the left keeps
  everything inside $[1,y]$:
  \[
    \Phi^h(v):=\int\Phi^\sharp(ve^\tau)h(\tau)\,d\tau=\frac1{2\pi i}\int_{(\eta_y)}\mathfrak G(s)\,\hat h_0(s/T)\,v^s\,\frac{ds}s,
    \qquad|\hat h_0(w)|\le C_me^{2|\mathrm{Re}\,w|}(1+|w|)^{-m}.
  \]
\item Shift the contour to $\mathrm{Re}\,s=-f\delta_{\mathrm{zf}}(T_1)$ for $|t|\le T_1\le T_y$. There is no residue at $s=0$,
  because $\mathfrak G(0)=0$. Take $T:=e^{\sqrt{\log v}}$, $T_1:=e^{2\sqrt{\log v}}(\log y)^{2/m}$ and $m:=19$.
\item The contributions are as follows. The shifted line gives at most
  $K_8\,v^{-f\delta_{\mathrm{zf}}(T_1)}\log^2(T_1+3)$. The horizontal segments, and the line $\mathrm{Re}\,s=\eta_y$ beyond
  $T_1$, give at most
  \[
    K_9v^{\eta_y}(\log\log y)(T/T_1)^m\le e^{-\sqrt{\log v}},
  \]
  using $v^{\eta_y}\le(\log y)^{17}$. The
  smoothing error is
  \[
    |\Phi^\sharp(v)-\Phi^h(v)|\le\int_{(ve^{-2/T},v]}\frac{d\Pi+d\Pi_0}u\le\frac6T+3v^{\theta'-1},
  \]
  using $d\Pi\le dN$.
\item The heights used lie below $T_v$ once $\log v\ge1.1\cdot10^4\Lambda^2(\log\log v)^2$; here we apply the argument with
  $y':=v$ in place of $y$, which is allowed since (LP-a) holds on $[1,v]$; and $y'=v$ satisfies \eqref{eq:lp-yrange}.
  For $\log v\ge(1000+4\log\Lambda)^2$ the offset in
  \eqref{eq:lp-zf} at most doubles the denominator. Together this gives (i) with $c_{\mathrm{LP}}=c_{\mathrm{zf}}/5$. The
  threshold $X_{\mathrm{LP}}$ is chosen above both heights, so $\log X_{\mathrm{LP}}=O(\Lambda^2\log^2\Lambda)$, and
  $\log K=O(\Lambda)$.
\item (ii) is the same argument with $\mathfrak G_1(s):=\log(\zeta^\sharp(s)(s-1)/s)$, the kernel
  $v^s\hat h_0((s-1)/T)/s$, and the contour near $\sigma=1$.
\end{itemize}

\emph{Step 10: short windows; proof of (iii).} With the smooth window $\varphi$, write $\int\varphi\,d(\Pi^\sharp-\Pi_0)$ as the
contour integral of $\log(\zeta^\sharp(s)(s-1)/s)$ against the Mellin transform $\tilde\varphi$ of $\varphi$. Since $\varphi$ has
transition layers $\eta\Lambda_p^{-2}$, $\int|\tilde\varphi(\sigma+it)|\,dt\le4.04\,C_\varphi v^\sigma\Lambda_p^2$ with
$C_\varphi\le4$, and the integrand is negligible beyond height $T_1=e^5\Lambda_p^2/\eta$, i.e.\ $L_1=\log T_1$. No $1/T$
smoothing term arises, so $T_1$ is fixed by $\eta$ alone. Shifting to $\sigma=1-f\delta_{\mathrm{zf}}(T_1)$ and using Step~7 for
the width at height $e^{L_1}$ gives \eqref{eq:lpplus}.
\begin{itemize}
\item \emph{Range.} Step~7 gives the width at height $e^{L_1}$ if \eqref{eq:lp-yrange} holds and $L_1\le\log T_y$. Both hold for
  $\Lambda_p\ge76\Lambda^3$. First, $\log y\ge\Lambda_p\ge76\Lambda^3$, and $\log y/\log\log y$ is increasing, so
  \eqref{eq:lp-yrange} holds. Second, $L_1+500\le2\Lambda+2\log\Lambda_p+505$, while $\log T_y+500\ge\Lambda_p/(51\Lambda\log\Lambda_p)$;
  the latter exceeds the former for every $\Lambda_p\ge76\Lambda^3$, since both sides increase and at $\Lambda_p=76\Lambda^3$ the
  right side is $\ge\Lambda^2/(3\log\Lambda+5)\gg2\Lambda$. Also $L_1\ge2\log\Lambda_p\ge40$. So $5\le L_1\le\log T_y$, and the
  bound on $G(L_1)$ of Step~8 applies.
\item \emph{Size.} For $\eta\ge e^{-2\Lambda}$ and $L_1\le e^{\Lambda/2}$ (so that $\log G(L_1)\le\Lambda$ by Step~8),
  $\mathrm{Rel}\le0.1$ needs $0.9c_{\mathrm{zf}}\Lambda_p/(L_1+442+2\log\Lambda+\log L_1)\ge3\Lambda+3\log\Lambda_p+6$, that is,
  $\Lambda_p\ge76\Lambda^3$. If $L_1>e^{\Lambda/2}$, then $\log\Lambda_p\ge e^{\Lambda/2}/3$ and $G(L_1)\le\Lambda_p$ by Step~8, and the
  same comparison holds with room to spare. The threshold $76\Lambda^3$ is below $\Lambda_1=100\Lambda^3$ of
  Lemma~\ref{lem:si}.\qedhere
\end{itemize}
\end{proof}

\begin{remark}\label{rem:lp}
Positivity cannot be dropped: $N=\lfloor u\rfloor+\frac23\cos(3\log u+\pi/2)$ has $|N-u|\le2$, but its $\zeta^\sharp$ has a zero
near $1+3i$. The theorem uses only the realised $N$ on $[1,y]$, so it holds pathwise on $\Good(<y)$, and the adaptedness of $T$
plays no role. Parts (i) and (ii) are used only in Proposition~\ref{prop:g1}(B), whose range starts at
$X_{\mathrm{sw}}\gg X_{\mathrm{LP}}$; part (iii) is used in Lemma~\ref{lem:si}.
\end{remark}
